\section{Time, conversion, and the two square roots of the pairing}
\label{sec:clock}

\subsection{One conversion constant}

The split form is not dimensionally innocent. If $V^+$ carries times and $V^-$ lengths, then
$\eta(v,v)=\sum_i(t^i)^2-\sum_a(x^a)^2$ adds squared times to squared lengths, and a conversion is needed, a priori one
constant for each temporal axis.

\begin{proposition}[A pairing alone does not give isotropy]
\label{prop:KC}
A map exchanging $V^+$ and $V^-$ cannot carry a single dimension; it converts times to lengths one way and lengths to times
the other. In block form
\begin{equation}
\label{eq:KC}
\cK_C=\begin{pmatrix}0&C^{-1}\\ C&0\end{pmatrix},\qquad C:V^+\to V^-,\qquad [C]=LT^{-1},
\end{equation}
one has $\cK_C^2=\id$ for every invertible $C$, including $C=\operatorname{diag}(c_1,c_2,c_3)$ with distinct entries. An
involutive pairing therefore does not force $c_1=c_2=c_3$.
\end{proposition}

\begin{theorem}[One conversion constant]
\label{thm:onec}
If a common $SO(3)$ acts on $V^+$ and on $V^-$ by the vector representation and $C$ intertwines the two actions, then
$C=c\,\id$ with $[c]=LT^{-1}$.
\end{theorem}

\begin{proof}
The vector representation of $SO(3)$ is irreducible of real type, so its commutant is $\R$ and Schur's lemma places $C$ in
$\R\,\id$ \cite{GoodmanWallach}.
\end{proof}

\subsection{The two square roots of the pairing}

Dimensional analysis fixes $[C]$ and leaves the sign of the lower-left block free. The two choices are different operators,
and the rest of the paper turns on the difference.

\begin{proposition}[The grading root and the binding root]
\label{prop:twoK}
In units with $c=1$, and in block form on $V^+\oplus V^-$, put
\[
\cK=\begin{pmatrix}0&-\id\\ -\id&0\end{pmatrix},\qquad
G=\begin{pmatrix}0&-\id\\ \id&0\end{pmatrix},
\]
that is $\cK e_i=-e_{i+3}$, $\cK e_{i+3}=-e_i$, and $Ge_i=e_{i+3}$, $Ge_{i+3}=-e_i$. Then:
\begin{center}
\begin{tabular}{@{}lcccl@{}}
\toprule
& square & $\eta$-adjoint & $\eta(X\cdot,X\cdot)$ & role\\
\midrule
$\cK$ & $+\id$ & $\cK^\dagger=-\cK$ & $-\eta$ & an element of $\so{3,3}$; grades\\
$G$ & $-\id$ & $G^\dagger=+G$ & $-\eta$ & a complex structure outside $\so{3,3}$; binds\\
\bottomrule
\end{tabular}
\end{center}
with $X^\dagger:=\eta^{-1}X^{\!\top}\eta$. Both are anti-isometries. The eigenspaces of $\cK$ are the two totally null three-planes
spanned by $e_i\mp e_{i+3}$.
\end{proposition}

Both roots can be built inside the algebra, and they are related by a single trivector.

\begin{lemma}[Both roots from Clifford data]
\label{lem:Gclifford}
Let $\mathsf P=\sum_i e_ie_{i+3}$ and $T=e_1e_2e_3$, and define $\cK(v)=\tfrac12[\mathsf P,v]$ and $S(v)=TvT^{-1}$ on vectors. Then
$\cK(e_i)=-e_{i+3}$, $\cK(e_{i+3})=-e_i$, so $\cK^2=\id$; $S$ fixes the temporal and negates the spatial generators, so $S=\eta$ as a
map of $V$; and
\[
G=-\cK\circ S .
\]
\end{lemma}

\begin{proof}
$[e_ie_{i+3},e_i]=e_ie_{i+3}e_i-e_i^2e_{i+3}=-2e_{i+3}$ and $[e_ie_{i+3},e_{i+3}]=e_ie_{i+3}^2-e_{i+3}e_ie_{i+3}=-2e_i$, while the
other pole operators commute with both by \eqref{eq:bladecomm}. $T$ is a product of the three temporal generators, so it
commutes with each of them up to the sign $(-1)^{3-1}=+1$ and anticommutes with each spatial generator. Finally
$-\cK(S e_i)=-\cK(e_i)=e_{i+3}=Ge_i$ and $-\cK(Se_{i+3})=\cK(e_{i+3})=-e_i=Ge_{i+3}$.
\end{proof}

So the root that grades is $\tfrac12\ad_{\mathsf P}$, the adjoint action of the sum of the pole operators, and the root that binds
is the same map composed with the Cartan involution $S=\Ad_{e_1e_2e_3}$. That the two roots are distinguished exactly by the
Cartan involution is the cleanest single statement the flat and the curved analyses below jointly support.

\paragraph{Shift operators.} A clock coordinate $\lambda$ has dimension $T$, the operator $G$ is dimensionless, and their sum
requires a scale. With a length $\ell$,
\begin{equation}
\label{eq:kappa}
\kappa=\frac c\ell,\qquad \alpha_\pm=\partial_\lambda\pm\kappa G .
\end{equation}
The pointwise $\eta$-adjoint of $\kappa G$ is $\kappa G$, while the formal adjoint of the differential operator, after an integration by
parts with compatible boundary conditions, is $(\partial_\lambda+\kappa G)^*=-\partial_\lambda+\kappa G^\dagger=-\alpha_-$. The two
adjoints must be kept apart: the factorisations below are insensitive to the overall sign, positivity statements are not.

\subsection{The clock}
\label{sec:clocksimplex}

A clock is a unit vector of $V^+$. Which one? Let $w=(w_1,w_2,w_3)$ be normalised non-negative weights on the three temporal
axes, so that $w$ ranges over the two-simplex $\Delta^2$. The map $x_i=2\sqrt{w_i}$ sends $\Delta^2$ onto the part of the sphere
$\sum_ix_i^2=4$ in the closed positive octant, and since $dx_i=dw_i/\sqrt{w_i}$ the Euclidean metric pulls back to
\[
\sum_idx_i^2=\sum_i\frac{dw_i^2}{w_i},
\]
the Fisher--Rao metric. That metric is not a modelling choice: by \v{C}encov's theorem it is, up to scale, the only metric on a
finite simplex invariant under Markov morphisms \cite{Cencov,AmariNagaoka}. The space of clocks is therefore an octant of a round
sphere of radius two.

\begin{proposition}[The symmetric clock]
\label{prop:clock}
The barycentre $w=(\tfrac13,\tfrac13,\tfrac13)$ is the unique point of $\Delta^2$ fixed by the permutations of the three axes, and the
unique maximum of the entropy $-\sum_iw_i\log w_i$, with value $\log3$ and Hessian eigenvalues $-9$ and $-3$ in the chart
$(w_1,w_2)$. Its amplitude direction is the unit vector
\[
\widehat W=\frac{e_1+e_2+e_3}{\sqrt3},\qquad \widehat W^2=\id,\qquad
\lambda:=\eta(\widehat W,v)=\frac{t^1+t^2+t^3}{\sqrt3}=\sqrt3\,\bar t .
\]
\end{proposition}

The permutation argument carries the weight here; the entropy argument is Jaynes' principle and is a hypothesis, recorded as
such. What the proposition says is modest: if a clock is to treat the three temporal axes alike, it is $\widehat W$, and the
physical time it reads is $\sqrt3$ times the mean of the three constituent times. It does not say that $\lambda$ is a proper time,
that it is invariant under a group, or that the directions orthogonal to $\widehat W$ are unphysical. Those directions,
$\widehat W^\perp\cap V^+=\{t^1+t^2+t^3=0\}$, spanned by $e_1-e_2$ and $e_1+e_2-2e_3$, are the \emph{relative times}.

\begin{figure}[t]
\centering
\begin{tikzpicture}[scale=1.15]
% simplex
\begin{scope}
\coordinate (A) at (0,0); \coordinate (B) at (3,0); \coordinate (Cc) at (1.5,2.6);
\fill[tcol!8] (A)--(B)--(Cc)--cycle; \draw[tcol,thick] (A)--(B)--(Cc)--cycle;
\node[below left] at (A) {$e_1$}; \node[below right] at (B) {$e_2$}; \node[above] at (Cc) {$e_3$};
\coordinate (bar) at (1.5,0.8667);
\fill[scol] (bar) circle (2.3pt);
\node at (1.5,0.2) {\footnotesize $w=(\tfrac13,\tfrac13,\tfrac13)$};
\foreach \r in {0.22,0.42}{\draw[gray!60,thin] (bar) circle (\r);}
\node at (1.5,-0.6) {\small simplex $\Delta^2$ of clock weights};
\end{scope}
\draw[->,thick] (3.5,1.2) -- node[above] {\small $x_i=2\sqrt{w_i}$} (5.3,1.2);
% octant of sphere, drawn in an oblique projection
\begin{scope}[shift={(7.3,0)}]
\coordinate (O) at (0,0);
\coordinate (X) at (2.6,-0.25); \coordinate (Y) at (-0.9,-0.9); \coordinate (Z) at (0,2.6);
\draw[->,gray] (O)--($1.25*(X)$) node[right] {\small $x_2$};
\draw[->,gray] (O)--($1.35*(Y)$) node[left] {\small $x_1$};
\draw[->,gray] (O)--($1.15*(Z)$) node[above] {\small $x_3$};
\fill[tcol!12] (X) .. controls ($(X)+(0,1.8)$) and ($(Z)+(1.2,0)$) .. (Z)
  .. controls ($(Z)+(-0.8,-0.6)$) and ($(Y)+(-0.5,1.3)$) .. (Y)
  .. controls ($(Y)+(1.1,-0.2)$) and ($(X)+(-1.2,-0.6)$) .. (X);
\draw[tcol,thick] (X) .. controls ($(X)+(0,1.8)$) and ($(Z)+(1.2,0)$) .. (Z);
\draw[tcol,thick] (Z) .. controls ($(Z)+(-0.8,-0.6)$) and ($(Y)+(-0.5,1.3)$) .. (Y);
\draw[tcol,thick] (Y) .. controls ($(Y)+(1.1,-0.2)$) and ($(X)+(-1.2,-0.6)$) .. (X);
\coordinate (Wp) at (0.72,0.95);
\fill[scol] (Wp) circle (2.3pt);
\draw[->,scol,thick] (O) -- (Wp) node[above right] {\small $\widehat W$};
\node[align=center,anchor=north] at (1.4,-1.45) {\small octant of the sphere $\sum x_i^2=4$\\[-1pt] \small (Fisher--Rao metric)};
\end{scope}
\end{tikzpicture}
\caption{The space of clocks. Normalised weights on the three temporal axes form a simplex; the amplitude map carries it onto an
octant of a round sphere, whose round metric is the Fisher--Rao metric. The barycentre, the only permutation-invariant point and
the entropy maximum, corresponds to the symmetric clock $\widehat W=(e_1+e_2+e_3)/\sqrt3$ (\cref{prop:clock}).}
\label{fig:clock}
\end{figure}

\subsection{The spectrum of a conversion operator}

An operator that converts times into lengths is a mixed bivector $B_C=\sum_{i,a}C_{ia}\,e_ie_{a+3}$, symmetric by
\cref{thm:sig}. Its square controls the factorised interior dynamics of \cref{sec:shifted}, where the equation of motion becomes
$\ddot X=\kappa^2B_C^2X$, so its spectrum is worth knowing exactly.

\begin{theorem}[Spectrum of a conversion operator]
\label{thm:spectrumBC}
Write $C=\sum_k\sigma_ku_kv_k^{\!\top}$ in singular-value form and put $U_k=\sum_iu_k^ie_i$, $V_k=\sum_av_k^ae_{a+3}$ and
$B_k=U_kV_k$. Then $B_k^2=\id$, the $B_k$ commute, and
\[
B_C^2=\|C\|_F^2\,\id+2\sum_{k<l}\sigma_k\sigma_l\,\Omega_{kl},\qquad \Omega_{kl}=B_kB_l,\qquad \Omega_{12}\Omega_{13}=\Omega_{23}.
\]
Consequently
\[
\spec\bigl(B_C^2\bigr)=\bigl\{(\eta_1\sigma_1+\eta_2\sigma_2+\eta_3\sigma_3)^2:\ \eta\in\{\pm1\}^3\bigr\},
\]
four values, each of multiplicity two on $\R^8$.
\end{theorem}

\begin{proof}
$U_k$ and $V_k$ are orthogonal unit vectors of square $+\id$ and $-\id$, so $B_k^2=-U_k^2V_k^2=\id$; for $k\neq l$ the four vectors
$U_k,V_k,U_l,V_l$ are mutually orthogonal and $B_k,B_l$ commute by \eqref{eq:bladecomm}. Expanding the square gives the formula. The
seven nontrivial products of the commuting involutions $B_k$ are, in a rotated orthonormal basis of $V$, blades of grade two, four
or six, hence traceless; the joint eigenspace with signs $\eta$ therefore has dimension $\tfrac18\Tr\prod_k(\id+\eta_kB_k)=1$, and on it
$B_C=\eta\cdot\sigma$.
\end{proof}

\begin{corollary}[Rank counts clocks]
\label{cor:rankclock}
$B_C^2$ is a multiple of $\id$ if and only if $\rank C\le1$, and some eigenvalue of $B_C^2$ vanishes if and only if one singular value
equals the sum of the other two. The rank of $C$ is the number of independent temporal directions the conversion uses.
\end{corollary}

\begin{example}[The symmetric clock is dynamically blind]
\label{ex:blindclock}
The operators $G_a:=\widehat W e_{a+3}$, which convert along the symmetric clock, have $C=\widehat W\hat e_a^{\!\top}$ of rank one and
singular values $(1,0,0)$. Hence $G_a^2=\id$, the factorised equation of motion is $\ddot X=\kappa^2X$, and the conversion operator does not
appear in it at all. The clock singled out by the symmetry is exactly the clock whose dynamics cannot see the Clifford structure. This is
forced by the rank: if every left singular vector of $C$ is proportional to $\widehat W$, then $C$ has rank one. A conversion visible in the
equations of motion therefore needs a temporal direction transverse to $\widehat W$, that is, a relative time. Raising the rank does not
produce oscillation either: the eigenvalues $(\eta\cdot\sigma)^2$ are real and non-negative at every rank, so higher rank makes the
instability anisotropic, with four rates $\kappa|\eta\cdot\sigma|$ instead of one, and can create flat directions, but never a bounded
oscillation. Bounded motion has to come from elsewhere, and \cref{sec:neutral} shows where.
\end{example}
